\chapter{A Concrete Naming Certificate for $\LayeredRelationGateUp$}
\label{ch:concrete-instances-layered-relation-gate-namecert}
\origin{ai}

The $\LayeredRelationGateUp$ packet realizes the gate surface of \autoref{def:layered-relation-gate} as a finite same-schema relation verdict. It sits after the relation certificate of \autoref{ch:visions-layered-relation-certificate}, the concrete certificate packet of \autoref{ch:concrete-instances-layeredrelationcert-namecert}, the force inter-history relation packet of \autoref{ch:concrete-instances-forceinterhistrelation-namecert}, and the axiom-refusal ledger of \autoref{ch:concrete-instances-axiom-refusal-ledger-namecert}. Those siblings supply relation-certificate rows, force-relation rows, and refusal-ledger rows; this packet binds their gate verdict, identity refusal, and same-schema consumer route as one NameCert surface.

\falsifiablePrediction{A gate packet that lacks a not-preserved row cannot export a same-schema consumer theorem.}

\independenceWitness{layeredrelationcert, forceinterhistrelation, groundcompilerrecognizerflow, axiomrefusalledger: the gate carrier has a verdict and identity-refusal row absent from the certificate-only, force-relation, recognizer-flow, and axiom-refusal packets, so no bijection can preserve all NameCert rows.}

\begin{definition}[LayeredRelationGate carrier]
\label{def:layered-relation-gate-carrier}
A $\LayeredRelationGateUp$ carrier is a finite $\BHist$ packet
$$
K=(A,B,L,P,N,F,G,H,C,\Pi)
$$
with the following displayed rows:
\begin{enumerate}
\item $A$ and $B$ are the two source-domain rows whose relation is being gated;
\item $L$ is the finite layer-list row pairing source layers, carrier families, classifier families, and layer-map entries;
\item $P$ is the preserved row, listing exactly the relation content that may be exported under same-schema language;
\item $N$ is the not-preserved row, listing object identity, empirical law, proof transfer, subject-substance, global-frame, and bridge-strength content refused by the gate;
\item $F$ is the finite ledger, exactness, and failure-boundary row, recording each gate question of \autoref{def:layered-relation-gate};
\item $G$ is the gate verdict row, whose accepted value permits same-schema relation export only when $A,B,L,P,N,F$ are all displayed;
\item $H$ records componentwise $\hsame$ transport for source, layer, preserved, not-preserved, refusal, verdict, continuation, and provenance rows;
\item $C$ records displayed $\Cont$ routes for source inspection, layer-list replay, preserved-row export, not-preserved-row inspection, refusal-ledger replay, and consumer handoff;
\item $\Pi$ is the $\Pkg$ and local $\NameCert_{\LayeredRelationGateUp}$ row over \autoref{ch:visions-layered-relation-certificate}, \autoref{ch:concrete-instances-layeredrelationcert-namecert}, \autoref{ch:concrete-instances-forceinterhistrelation-namecert}, \autoref{ch:concrete-instances-groundcompilerrecognizerflow-namecert}, \autoref{ch:concrete-instances-axiom-refusal-ledger-namecert}, $\BHist$, $\hsame$, $\Cont$, $\Pkg$, and $\NameCert$.
\end{enumerate}
The carrier has no coordinate for identifying $A$ with $B$ as objects. Same-schema relation language is an export of the verdict row $G$ after the preserved row $P$ and the not-preserved row $N$ have both been displayed.
\end{definition}

\begin{theorem}[LayeredRelationGate NameCert obligations]
\label{thm:layered-relation-gate-namecert-obligations}
For the carrier of \autoref{def:layered-relation-gate-carrier}, the local $\NameCert_{\LayeredRelationGateUp}$ surface consists of seven finite obligations:
\begin{enumerate}
\item carrier exposure for $K=(A,B,L,P,N,F,G,H,C,\Pi)$;
\item source-domain exposure for $A$ and $B$ before any same-schema relation read;
\item layer-list exposure for the paired source layers, carrier families, classifier families, and layer map in $L$;
\item preserved-row disclosure through $P$;
\item not-preserved disclosure through $N$ before verdict export;
\item refusal-ledger and failure-boundary replay through $F$;
\item consumer handoff through $G,H,C,\Pi$ without object identity, empirical law, proof transfer, subject-substance, global-frame, bridge-strength, or hidden-choice export.
\end{enumerate}
\end{theorem}
\begin{proof}
Carrier exposure is projection from the finite tuple in \autoref{def:layered-relation-gate-carrier}. The accepted packet has no fields outside $A,B,L,P,N,F,G,H,C,\Pi$.

Source-domain and layer-list exposure are ordered by the displayed continuation row $C$. A same-schema consumer first reads the source rows $A,B$ and then the layer-list row $L$ before any preserved or refused content is available.

The preserved and not-preserved obligations are the rows $P$ and $N$. Since $G$ can be read only after both rows have been displayed, a passing gate cannot hide refused content behind same-schema wording.

The refusal-ledger obligation is the row $F$. It records the failure boundary for the gate questions of \autoref{def:layered-relation-gate}, including identity collapse, hidden empirical law, subject-substance, global-frame, object-level $\mathsf{T}$, and strength over-reading.

Finally, the structural rows $H,C,\Pi$ transport, replay, package, and name the same finite verdict surface. They do not add object identity or choice rows, so the local NameCert surface is exactly the listed gate packet.
\end{proof}

\begin{theorem}[LayeredRelationGate same-schema non-identity]
\label{thm:layered-relation-gate-same-schema-nonidentity}
If an accepted $\LayeredRelationGateUp$ packet passes its verdict row $G$, then downstream consumers may use the displayed same-schema relation over $A,B,L,P,N,F$. They may not read the verdict as object identity between $A$ and $B$, as empirical law, as proof transfer, as subject-substance transfer, as a global-frame theorem, or as bridge strength.
\end{theorem}
\begin{proof}
Begin with the NameCert obligations of \autoref{thm:layered-relation-gate-namecert-obligations}. The verdict is downstream from the source rows, layer-list row, preserved row, not-preserved row, and refusal ledger.

The same-schema relation is therefore exhausted by $A,B,L,P,N,F,G$. The preserved row $P$ records what may be exported, while the not-preserved row $N$ records what the export refuses.

Object identity and the other excluded readings are listed only in the not-preserved and refusal rows. Since neither $H$ nor $C$ nor $\Pi$ contains a constructor from those rows to preserved content, a passing verdict licenses relation use without promoting the two domains to one object.
\end{proof}

\begin{theorem}[LayeredRelationGate consumer route]
\label{thm:layered-relation-gate-consumer-route}
Every same-schema consumer of a passing $\LayeredRelationGateUp$ packet factors through the relation-certificate source of \autoref{ch:visions-layered-relation-certificate}, the concrete certificate surface of \autoref{ch:concrete-instances-layeredrelationcert-namecert}, the force inter-history relation surface of \autoref{ch:concrete-instances-forceinterhistrelation-namecert}, and the axiom-refusal ledger of \autoref{ch:concrete-instances-axiom-refusal-ledger-namecert}. The consumer reads exactly the source, layer-list, preserved, not-preserved, ledger, exactness, failure-boundary, verdict, transport, continuation, package, and local NameCert rows of \autoref{def:layered-relation-gate-carrier}.
\end{theorem}
\begin{proof}
First read the relation-certificate discipline of \autoref{ch:visions-layered-relation-certificate} and \autoref{ch:concrete-instances-layeredrelationcert-namecert}. These supply the source, layer, preserved, not-preserved, exactness, and failure-boundary rows consumed by the gate.

Next use \autoref{ch:concrete-instances-forceinterhistrelation-namecert}. It supplies a concrete same-schema relation surface without turning the compared histories or domains into one object.

Now use \autoref{ch:concrete-instances-axiom-refusal-ledger-namecert}. It supplies the refusal-ledger discipline that keeps hidden axioms, global-frame imports, and unrecorded strength claims outside the exported relation.

Finally apply \autoref{thm:layered-relation-gate-namecert-obligations} and \autoref{thm:layered-relation-gate-same-schema-nonidentity}. They force the consumer through $A,B,L,P,N,F,G,H,C,\Pi$ and make the verdict a same-schema relation gate, not an identity theorem.
\end{proof}

The Lean follow-up must provide $\Nontrivial\ \LayeredRelationGateUp$ and $\FieldFaithful\ \LayeredRelationGateUp$ instances in \texttt{lean4/BEDC/Derived/LayeredRelationGateUp/TasteGate.lean}.

\closureat{LayeredRelationGateUp}{seedStr}

\begin{closurestatus}{\LayeredRelationGateUp}
  \constructivestory{A $\LayeredRelationGateUp$ packet is generated from source domains, finite layer lists, preserved and not-preserved rows, a refusal ledger, a verdict row, componentwise $\hsame$ transport, displayed $\Cont$ routes, and local $\Pkg$ and $\NameCert$ rows.}
  \theoryclosure{\seedClosure}
  \scopeclosed{\autoref{def:layered-relation-gate-carrier}, \autoref{thm:layered-relation-gate-namecert-obligations}, \autoref{thm:layered-relation-gate-same-schema-nonidentity}, and \autoref{thm:layered-relation-gate-consumer-route} bind the finite LayeredRelationGate carrier, NameCert obligations, same-schema non-identity verdict, and consumer route.}
  \formalstatus{\unformalizedV}
  \bridgestatus{none}
  \origin{ai}
  \notclaimed{This chapter does not close object identity between compared domains, empirical physics, proof transfer, subject-substance content, global-frame structure, Lean verification, audit cleanliness, or bridge checking.}
  \upgradepath{Obligation closure requires checked carrier admission, verdict determinacy, preserved and not-preserved disclosure, refusal-ledger replay, consumer handoff, and TasteGate field faithfulness for the same finite packet.}
\end{closurestatus}
